%%% =======================================================================
%%% ISLS Extension Phase 2 v1.0.0
%%% C16: Topological-Spectral Orbit Core (kernel-k13 Absorption)
%%% C17: Structured Persistence Layer (merkaba-store Absorption)
%%% Author: Sebastian Klemm
%%% Date: 16 March 2026
%%% Status: Implementation Specification for Claude Code
%%% =======================================================================
\documentclass[11pt,a4paper,twoside]{article}

\usepackage[utf8]{inputenc}
\usepackage{amsmath,amssymb,amsthm,mathtools}
\usepackage{algorithm,algorithmic}
\usepackage{booktabs,longtable,tabularx}
\usepackage{enumitem}
\usepackage{geometry}
\usepackage{hyperref}
\usepackage{cleveref}
\usepackage{xcolor}
\usepackage{fancyhdr}

\geometry{margin=2.5cm,top=3cm,bottom=3cm,headheight=14pt}

\theoremstyle{definition}
\newtheorem{definition}{Definition}[section]
\newtheorem{axiom}{Axiom}[section]
\newtheorem{requirement}{Requirement}[section]
\newtheorem{invariant}{Invariant}[section]

\theoremstyle{plain}
\newtheorem{theorem}{Theorem}[section]
\newtheorem{proposition}{Proposition}[section]
\newtheorem{lemma}{Lemma}[section]

\theoremstyle{remark}
\newtheorem{remark}{Remark}[section]

\newcommand{\ISLS}{\textsf{ISLS}}
\newcommand{\RR}{\mathbb{R}}
\newcommand{\CC}{\mathbb{C}}
\newcommand{\NN}{\mathbb{N}}
\newcommand{\ZZ}{\mathbb{Z}}
\newcommand{\Hfive}{H^{5}}
\newcommand{\calL}{\mathcal{L}}
\newcommand{\calG}{\mathcal{G}}
\newcommand{\calR}{\mathcal{R}}
\newcommand{\calS}{\mathcal{S}}
\newcommand{\calK}{\mathcal{K}}
\newcommand{\calP}{\mathcal{P}}
\newcommand{\calT}{\mathcal{T}}
\DeclareMathOperator{\Tr}{Tr}
\DeclareMathOperator{\diag}{diag}
\DeclareMathOperator{\spec}{spec}
\DeclareMathOperator{\Dig}{Dig}
\DeclareMathOperator{\Canon}{Canon}
\newcommand{\norm}[1]{\left\lVert #1 \right\rVert}
\newcommand{\abs}[1]{\left\lvert #1 \right\rvert}
\newcommand{\innerprod}[2]{\langle #1, #2 \rangle}

\pagestyle{fancy}
\fancyhf{}
\fancyhead[LE,RO]{\ISLS{} Extension Phase 2 v1.0.0}
\fancyhead[RE,LO]{Sebastian Klemm}
\fancyfoot[C]{\thepage}

\hypersetup{
  colorlinks=true,
  linkcolor=blue!60!black,
  pdftitle={ISLS Extension Phase 2 v1.0.0},
  pdfauthor={Sebastian Klemm},
}

\begin{document}

\begin{titlepage}
\centering
\vspace*{2.5cm}
{\Huge\bfseries \ISLS{} Extension Phase~2\\[0.3cm]
v1.0.0}\\[1.5cm]
{\Large C16: Topological-Spectral Orbit Core\\
C17: Structured Persistence Layer}\\[1cm]
{\large Absorbing \texttt{kernel-k13} and \texttt{merkaba-store}\\
into the Intelligent Semantic Ledger Substrate}\\[2cm]
{\large Sebastian Klemm}\\[0.5cm]
{\large 16 March 2026}\\[2cm]

\begin{abstract}
\noindent
This document specifies two additional extension crates for \ISLS{}:
(1)~\texttt{isls-topology} (C16), absorbing the topological-spectral
machinery from \texttt{kernel-k13} of the MERKABA Factory workspace,
providing spectral graph analysis via the graph Laplacian, signal
propagation via continuous-time quantum walks (CTQW), and phase
synchronization via Kuramoto coupling; and
(2)~\texttt{isls-store} (C17), absorbing the structured SQLite
persistence model from \texttt{merkaba-store}, upgrading \ISLS{} from
filesystem-based storage to a relational persistence layer with
projects, runs, crystals, settings, and maintenance operations.

C16 constitutes the \emph{orbit core}: it transforms the persistent
graph from a static data structure into a dynamic spectral-topological
system whose invariants feed directly into constraint extraction,
gate evaluation, and crystal topological signatures. C17 constitutes
the \emph{product shell}: it makes \ISLS{} operable as a persistent
workspace rather than a stateless pipeline.

All extensions preserve invariants I1--I20 and all existing tests.

\medskip\noindent\textbf{Status:} Implementation Specification.\\
\textbf{Parent:} ISLS v1.0.0 + Extension Architecture v1.0.0.\\
\textbf{Absorbs:} \texttt{kernel-k13}, \texttt{merkaba-store}.
\end{abstract}
\end{titlepage}

\tableofcontents
\newpage

%%% =====================================================================
\part{Foundations}

\section{Scope and Absorption Policy}

\subsection{What Is Absorbed}

From \texttt{kernel-k13}:
\begin{itemize}[nosep]
  \item Spectral graph analysis (Laplacian, eigenvalues, spectral gap).
  \item Continuous-time quantum walk (CTQW) signal propagation.
  \item Kuramoto phase synchronization dynamics.
  \item Discrete topology logic (DTL) predicates.
\end{itemize}

From \texttt{merkaba-store}:
\begin{itemize}[nosep]
  \item SQLite-backed persistence with typed entities.
  \item Project, run, crystal, settings data model.
  \item Migration framework.
  \item Maintenance operations (vacuum, integrity check).
\end{itemize}

\subsection{What Is Not Absorbed}

\begin{itemize}[nosep]
  \item \texttt{kernel-qmh} (quantum-hybrid proxy): too speculative for
        the domain-agnostic core. Remains available as an optional
        future extension.
  \item \texttt{domain-translator}: language-specific frontends are out
        of scope for \ISLS{} core.
  \item \texttt{domain-artifact}: DAE synthesis is a domain application,
        not substrate infrastructure.
  \item \texttt{merkaba-app}: GUI is an application layer, not a
        substrate component. May be built later on top of C17.
\end{itemize}

\subsection{Absorption Criterion}

A component from MERKABA Factory is absorbed if and only if:
\begin{enumerate}[nosep,label=(\roman*)]
  \item it is domain-agnostic (no language-specific or UI-specific logic),
  \item it strengthens at least one existing \ISLS{} invariant or metric,
  \item it can be expressed as a pure function of the persistent graph
        and configuration (no external state dependencies),
  \item it preserves determinism under fixed $(\Sigma, \mathrm{RD})$.
\end{enumerate}

\section{Invariant Preservation}

\begin{requirement}[Phase 2 Non-Regression]\label{req:p2nonreg}
All extensions in this document \textbf{MUST} preserve invariants
I1--I20, all acceptance tests AT-01 through AT-20, all extension tests
AT-R1 through AT-S5, and all 161+ existing tests. The final test count
after Phase~2 \textbf{MUST} be $\ge 181$.
\end{requirement}


%%% =====================================================================
\part{Crate C16: \texttt{isls-topology}}\label{part:topology}

\section{Purpose}\label{sec:topo:purpose}

The topology crate transforms the persistent graph from a static
adjacency structure into a dynamic spectral-topological system. It
computes three families of invariants that feed into constraint
extraction, gate evaluation, crystal signatures, and the spiral
scheduler:

\begin{enumerate}[nosep,label=(\alph*)]
  \item \textbf{Spectral invariants}: eigenvalues and eigenvectors of
        the graph Laplacian, spectral gap, algebraic connectivity,
        Cheeger constant estimate.
  \item \textbf{Propagation invariants}: signal spreading speed and
        reachability via continuous-time quantum walk (CTQW) amplitudes.
  \item \textbf{Synchronization invariants}: Kuramoto order parameter
        measuring phase coherence of entity clusters.
\end{enumerate}

Together these form the \emph{orbit core}: a continuously updated
spectral portrait of the graph that serves as the dynamical backbone
for all higher layers.


\section{Mathematical Foundations}\label{sec:topo:math}

\subsection{Graph Laplacian}

\begin{definition}[Weighted Adjacency Matrix]\label{def:adjmatrix}
Given the persistent graph $G = (V, E, \Sigma_E, \Phi, D)$ with
$|V| = n$ vertices, define the weighted adjacency matrix
$W \in \RR^{n \times n}$ by
\[
  W_{ij} =
  \begin{cases}
    w_e & \text{if edge } e = (v_i, v_j) \in E, \\
    0 & \text{otherwise},
  \end{cases}
\]
where $w_e$ is the resonance weight from the edge annotation
$\Sigma_E(e)$. For undirected analysis, symmetrize:
$\bar{W} = (W + W^\top) / 2$.
\end{definition}

\begin{definition}[Degree Matrix]\label{def:degmatrix}
The degree matrix $D \in \RR^{n \times n}$ is diagonal with
\[
  D_{ii} = \sum_{j=1}^{n} \bar{W}_{ij}.
\]
\end{definition}

\begin{definition}[Graph Laplacian]\label{def:laplacian}
The unnormalized graph Laplacian is
\[
  L = D - \bar{W}.
\]
The normalized graph Laplacian is
\[
  L_{\mathrm{norm}} = I - D^{-1/2} \bar{W}\, D^{-1/2},
\]
with convention $D^{-1/2}_{ii} = 0$ when $D_{ii} = 0$.
\end{definition}

\begin{definition}[Spectral Decomposition]\label{def:spectral}
Let $0 = \lambda_0 \le \lambda_1 \le \cdots \le \lambda_{n-1}$ be the
eigenvalues of $L$ with corresponding orthonormal eigenvectors
$u_0, u_1, \ldots, u_{n-1}$. The \emph{spectral gap} is
\[
  \Delta = \lambda_1,
\]
which measures algebraic connectivity. The \emph{Fiedler vector} is
$u_1$, the eigenvector corresponding to $\lambda_1$.
\end{definition}

\begin{proposition}[Spectral Gap and Connectivity]\label{prop:specgap}
$\Delta > 0$ if and only if $G$ is connected. A larger $\Delta$
indicates stronger connectivity and faster mixing.
\end{proposition}

\begin{proof}
Standard result from spectral graph theory. $\lambda_0 = 0$ with
eigenvector $\mathbf{1}$ always. The multiplicity of eigenvalue $0$
equals the number of connected components.
\end{proof}

\begin{definition}[Cheeger Constant Estimate]\label{def:cheeger}
The Cheeger constant $h(G)$ is bounded by the spectral gap via the
discrete Cheeger inequality:
\[
  \frac{\Delta}{2} \le h(G) \le \sqrt{2\Delta}.
\]
We use $\hat{h} = \sqrt{2\Delta}$ as the upper-bound estimate (computable
without solving the NP-hard exact Cheeger problem).
\end{definition}


\subsection{Continuous-Time Quantum Walk (CTQW)}

\begin{definition}[CTQW Propagator]\label{def:ctqw}
The CTQW propagator on the graph Laplacian is the unitary (for
Hermitian $L$) matrix exponential
\[
  U(t) = e^{-iLt}.
\]
For a signal initialized at vertex $j$ (basis state $|j\rangle$),
the amplitude at vertex $k$ after time $t$ is
\[
  \alpha_{jk}(t) = \langle k | U(t) | j \rangle
  = \sum_{m=0}^{n-1} e^{-i\lambda_m t}\, u_m(j)\, u_m(k).
\]
The transfer probability is $p_{jk}(t) = |\alpha_{jk}(t)|^2$.
\end{definition}

\begin{remark}
In the \ISLS{} context, CTQW is not a physical quantum process. It is
a \emph{classical computation} using the mathematical formalism of
quantum walks to measure how efficiently information propagates through
the graph. The ``quantum'' label refers to the mathematical structure,
not to hardware.
\end{remark}

\begin{definition}[Propagation Speed]\label{def:propspeed}
The propagation speed from vertex $j$ to vertex $k$ is defined as
\[
  v_{jk} = \frac{d_G(j,k)}{t^*_{jk}},
\]
where $d_G(j,k)$ is the shortest-path distance in $G$ and
$t^*_{jk} = \arg\min_t \{t > 0 : p_{jk}(t) \ge \theta_{\mathrm{prop}}\}$
is the first time the transfer probability exceeds threshold
$\theta_{\mathrm{prop}}$.
\end{definition}

\begin{definition}[Mean Propagation Time]\label{def:meanprop}
The mean propagation time of the graph is
\[
  \bar{t}_{\mathrm{prop}} = \frac{1}{|V|^2} \sum_{j,k \in V} t^*_{jk}.
\]
\end{definition}

\begin{definition}[CTQW Approximation]\label{def:ctqwapprox}
For computational efficiency, the CTQW propagator is approximated using
the spectral decomposition. Let $K \le n$ be the number of retained
eigenvalues (truncation rank). Then
\[
  U_K(t) = \sum_{m=0}^{K-1} e^{-i\lambda_m t}\, u_m u_m^\top.
\]
For $K = \min(n, K_{\max})$ with configurable $K_{\max}$, this
achieves $O(K^2 n)$ per time step instead of $O(n^3)$ for full
matrix exponentiation.
\end{definition}


\subsection{Kuramoto Synchronization}

\begin{definition}[Phase Assignment]\label{def:phaseassign}
Each vertex $v_i$ carries a phase $\phi_i \in [0, 2\pi)$ derived from
its 5D embedding:
\[
  \phi_i = \omega(\Phi(v_i)) \bmod 2\pi,
\]
where $\omega(\cdot)$ extracts the frequency coordinate from the
\texttt{FiveDState}.
\end{definition}

\begin{definition}[Kuramoto Dynamics]\label{def:kuramoto}
The Kuramoto model on the weighted graph evolves phases as
\[
  \dot{\phi}_i = \omega_i^{\mathrm{nat}} + \frac{\kappa}{D_{ii}}
  \sum_{j: (i,j) \in E} \bar{W}_{ij} \sin(\phi_j - \phi_i),
\]
where $\omega_i^{\mathrm{nat}}$ is the natural frequency of vertex $i$
(derived from embedding) and $\kappa > 0$ is the global coupling
strength.
\end{definition}

\begin{definition}[Kuramoto Order Parameter]\label{def:orderparameter}
The complex order parameter is
\[
  r\, e^{i\Psi} = \frac{1}{n} \sum_{j=1}^{n} e^{i\phi_j}.
\]
The magnitude $r \in [0,1]$ measures global phase coherence:
$r = 1$ means perfect synchronization, $r = 0$ means complete
incoherence. We define
\[
  r_k = |r\, e^{i\Psi}|
\]
at macro-step $k$.
\end{definition}

\begin{proposition}[Kuramoto Coherence and Mandorla]\label{prop:kuramandorla}
The Kuramoto order parameter $r_k$ provides a graph-intrinsic measure
of phase coherence that is independent of the carrier geometry. When
$r_k$ is high, the entities in the graph are phase-synchronized, which
corresponds to a physically meaningful mandorla coherence. Therefore
$r_k$ can serve as an alternative or complementary computation of the
coherence metric $Q_k$.
\end{proposition}

\begin{proof}[Proof sketch]
The mandorla coupling $\kappa(t)$ in the carrier geometry measures
interference between two helices. The Kuramoto $r_k$ measures
interference among all entity phases. Both converge to the same
qualitative judgment: high values indicate structural coherence,
low values indicate incoherence. The carrier-geometric $Q_k$ and the
graph-intrinsic $r_k$ are complementary views of the same phenomenon.
\end{proof}


\subsection{Discrete Topology Logic (DTL)}

\begin{definition}[DTL Predicates]\label{def:dtl}
A DTL predicate is a Boolean function on graph structure. The minimal
set of DTL predicates is:
\begin{enumerate}[nosep,label=(\roman*)]
  \item $\mathrm{Connected}(G)$: the graph has exactly one connected
        component ($\Delta > 0$).
  \item $\mathrm{Bipartite}(G)$: the graph is bipartite
        ($\lambda_{n-1} = 2$ for $L_{\mathrm{norm}}$, or all cycles
        have even length).
  \item $\mathrm{TreeLike}(G)$: $|E| = |V| - 1$ and $G$ is connected.
  \item $\mathrm{ClusterCount}(G, k)$: spectral clustering with $k$
        clusters produces non-empty clusters (multiplicity of
        near-zero eigenvalues $\approx k$).
  \item $\mathrm{Expander}(G, c)$: $\Delta \ge c$ (the graph is a
        $c$-spectral expander).
\end{enumerate}
\end{definition}

\begin{definition}[Topological Signature (Hardened)]\label{def:toposig}
The topological signature $\tau(C)$ of a crystal $C$ with condensed
region $G^*$ is extended to include:
\[
  \tau(C) = \langle \beta_0, \beta_1, \beta_2, \Delta, \hat{h},
            r_{\mathrm{sync}}, \bar{t}_{\mathrm{prop}},
            \mathrm{DTL}_{\mathrm{vec}} \rangle
\]
where:
\begin{itemize}[nosep]
  \item $\beta_0, \beta_1, \beta_2$: Betti numbers (existing),
  \item $\Delta$: spectral gap of the condensed region,
  \item $\hat{h}$: Cheeger constant estimate,
  \item $r_{\mathrm{sync}}$: Kuramoto order parameter of the region,
  \item $\bar{t}_{\mathrm{prop}}$: mean CTQW propagation time,
  \item $\mathrm{DTL}_{\mathrm{vec}}$: vector of DTL predicate results.
\end{itemize}
\end{definition}


\section{Integration with Existing Layers}\label{sec:topo:integration}

\subsection{Feeding Constraint Extraction (L2)}

\begin{requirement}[Spectral Enrichment of Extraction]\label{req:specextract}
Before \texttt{inverse\_weave} runs, the topology crate \textbf{MUST}
compute the spectral decomposition of the current graph and attach
the spectral gap $\Delta$, Fiedler vector $u_1$, and Kuramoto $r_k$
to the extraction context. Constraint candidates \textbf{MAY} use
these as additional features for coverage computation.
\end{requirement}

\begin{definition}[Spectral Constraint Templates]\label{def:specconstraint}
Two new constraint template families are added:
\begin{enumerate}[nosep,label=(\roman*)]
  \item \textbf{SpectralGap}: $\Delta \ge \Delta_{\min}$. The graph
        maintains minimum connectivity.
  \item \textbf{SyncCluster}: $r_k(\mathrm{cluster}_i) \ge r_{\min}$.
        A cluster of entities maintains minimum phase synchronization.
\end{enumerate}
\end{definition}

\subsection{Feeding Gate Evaluation (L3)}

\begin{definition}[New Metrics M25--M27]\label{def:newmetrics}
Three new metrics join the existing M1--M24:
\begin{itemize}[nosep]
  \item \textbf{M25 Spectral Gap}: $\Delta_k$ of the current graph.
        Alert threshold: $\Delta < 0.01$ (disconnection risk).
  \item \textbf{M26 Kuramoto Coherence}: $r_k$ of the full graph.
        Alert threshold: $r_k < 0.1$ (no synchronization).
  \item \textbf{M27 Mean Propagation Time}: $\bar{t}_{\mathrm{prop},k}$.
        Alert threshold: $\bar{t}_{\mathrm{prop}} > 100$ (signals
        propagate too slowly).
\end{itemize}
\end{definition}

\begin{remark}
M25--M27 are \emph{informational metrics}, not gate variables. They do
not add new gates to the Kairos condition. They enrich the crystal
topological signature and the metric dashboard. A future profile
\textbf{MAY} promote them to gate variables if desired.
\end{remark}

\subsection{Feeding the Spiral Scheduler (C15)}

\begin{requirement}[Topology-Informed Scheduling]\label{req:toposched}
The spiral scheduler \textbf{MAY} use spectral gap and Kuramoto
coherence as additional inputs to the scheduling function:
\[
  \sigma(d, F, S, \Delta, r_k, \theta_{\mathrm{sched}}).
\]
When $\Delta$ drops sharply (structural disconnection event) or $r_k$
drops sharply (desynchronization event), the scheduler \textbf{SHOULD}
increase sub-step resolution.
\end{requirement}

\subsection{Feeding Crystal Signatures}

\begin{requirement}[Hardened Topological Signature]\label{req:hardenedtopo}
Every crystal produced after C16 is active \textbf{MUST} include the
hardened topological signature $\tau(C)$ as defined in
\cref{def:toposig}. This extends the existing topology\_signature field
in the crystal schema.
\end{requirement}


\section{Fixpoint Detection}\label{sec:topo:fixpoint}

Absorbed from \texttt{kernel-engine}'s \texttt{FixpointDetector}.

\begin{definition}[Jaccard Fixpoint Detection]\label{def:jaccardfix}
Let $V_k \subseteq V$ be the set of ``active'' vertex identifiers at
macro-step $k$ (defined as vertices with at least one updated embedding
in the current tick). The Jaccard distance between consecutive steps is
\[
  d_J(V_k, V_{k-1}) = 1 - \frac{|V_k \cap V_{k-1}|}{|V_k \cup V_{k-1}|}.
\]
The system has reached a structural fixpoint when
$d_J(V_k, V_{k-1}) \le \epsilon_{\mathrm{fix}}$ for
$n_{\mathrm{consec}}$ consecutive steps.
\end{definition}

\begin{requirement}[Fixpoint Notification]\label{req:fixnotify}
When fixpoint is detected, the topology crate \textbf{MUST}:
\begin{enumerate}[nosep,label=(\alph*)]
  \item emit a \texttt{FixpointReached} event to the audit log,
  \item notify the spiral scheduler (which \textbf{MAY} reduce
        resolution to $n_{\min}$),
  \item \textbf{NOT} halt the engine (fixpoints can be transient; the
        engine must continue processing new observations).
\end{enumerate}
\end{requirement}


\section{Deduplication}\label{sec:topo:dedup}

Absorbed from the \texttt{vacuum-pipeline} canonicalization stage.

\begin{definition}[Observation Deduplication]\label{def:dedup}
Two observation records $o_1, o_2$ are duplicates if and only if
$\Dig(\Canon(o_1)) = \Dig(\Canon(o_2))$. The deduplication filter
is applied after canonicalization (L0) and before persistence (L1).
Duplicate observations are logged but not ingested.
\end{definition}

\begin{requirement}[Dedup Logging]\label{req:deduplog}
Every suppressed duplicate \textbf{MUST} be recorded in the audit log
with: original digest, duplicate count, source identifier, timestamp.
This preserves the evidence that duplicates were observed without
polluting the graph.
\end{requirement}


\section{Computational Complexity}\label{sec:topo:complexity}

\begin{tabularx}{\textwidth}{l l X}
\toprule
\textbf{Operation} & \textbf{Complexity} & \textbf{Notes} \\
\midrule
Laplacian construction & $O(|E|)$ & Sparse matrix from edge list \\
Spectral decomposition & $O(K^2 n)$ & $K$-truncated eigendecomp via
  Lanczos iteration; $K \le 50$ typical \\
CTQW propagation & $O(K^2 n \cdot T)$ & $T$ time steps, truncated
  spectral method \\
Kuramoto one step & $O(|E|)$ & Sparse neighbor sum \\
Kuramoto $N_{\mathrm{steps}}$ & $O(N_{\mathrm{steps}} \cdot |E|)$ &
  RK4 integration \\
Fixpoint detection & $O(|V|)$ & Jaccard on vertex sets \\
Deduplication & $O(1)$ amortized & Hash lookup in \texttt{BTreeSet} \\
\bottomrule
\end{tabularx}

\begin{requirement}[Topology Budget]\label{req:topobudget}
A conformant implementation \textbf{MUST} declare a per-tick budget
for topology computation. If the budget is exceeded (e.g., graph too
large for full spectral decomposition), the implementation \textbf{MUST}
fall back to:
\begin{enumerate}[nosep,label=(\alph*)]
  \item reduced truncation rank $K$,
  \item or sampling-based Kuramoto on a subgraph,
  \item or skipping CTQW propagation entirely (marking
        $\bar{t}_{\mathrm{prop}}$ as \texttt{N/A} in the signature).
\end{enumerate}
The fallback decision \textbf{MUST} be logged.
\end{requirement}


\section{Rust API}\label{sec:topo:api}

\begin{verbatim}
// isls-topology/src/lib.rs

pub struct SpectralDecomposition {
    pub eigenvalues: Vec<f64>,           // sorted ascending
    pub eigenvectors: Vec<Vec<f64>>,     // column-major
    pub spectral_gap: f64,              // lambda_1
    pub cheeger_estimate: f64,          // sqrt(2 * lambda_1)
    pub fiedler_vector: Vec<f64>,       // u_1
    pub truncation_rank: usize,         // K
}

pub struct CtqwResult {
    pub transfer_probabilities: Vec<Vec<f64>>,  // p[j][k] at time t*
    pub propagation_speeds: Vec<Vec<f64>>,       // v[j][k]
    pub mean_propagation_time: f64,              // t_bar
}

pub struct KuramotoState {
    pub phases: Vec<f64>,               // phi_i for each vertex
    pub order_parameter: f64,           // r (magnitude)
    pub mean_phase: f64,                // Psi
    pub cluster_coherences: Vec<f64>,   // r per spectral cluster
}

pub struct TopologicalSignature {
    pub betti_numbers: Vec<u64>,        // beta_0, beta_1, beta_2
    pub spectral_gap: f64,
    pub cheeger_estimate: f64,
    pub kuramoto_coherence: f64,
    pub mean_propagation_time: f64,
    pub dtl_predicates: BTreeMap<String, bool>,
}

pub struct FixpointDetector {
    epsilon: f64,
    n_consecutive: usize,
    max_iterations: u64,
    // internal state...
}

pub enum FixpointStatus {
    Converged { iteration: u64 },
    NotYet { jaccard_distance: f64 },
    MaxIterations,
}

pub struct TopologyConfig {
    pub spectral_k_max: usize,          // max eigenvalues to compute
    pub ctqw_time_max: f64,             // max propagation time
    pub ctqw_threshold: f64,            // theta_prop
    pub ctqw_time_steps: usize,         // discretization steps
    pub kuramoto_coupling: f64,         // kappa
    pub kuramoto_steps: usize,          // integration steps per tick
    pub kuramoto_dt: f64,               // RK4 step size
    pub fixpoint_epsilon: f64,          // Jaccard convergence threshold
    pub fixpoint_consecutive: usize,    // required consecutive stable
    pub budget_ms: u64,                 // per-tick compute budget
}

// Core functions
pub fn compute_laplacian(graph: &PersistentGraph) -> SparseLaplacian;
pub fn spectral_decompose(
    laplacian: &SparseLaplacian,
    k_max: usize,
) -> SpectralDecomposition;

pub fn ctqw_propagate(
    spectral: &SpectralDecomposition,
    config: &TopologyConfig,
) -> CtqwResult;

pub fn kuramoto_step(
    state: &mut KuramotoState,
    graph: &PersistentGraph,
    config: &TopologyConfig,
);

pub fn kuramoto_order_parameter(phases: &[f64]) -> (f64, f64);

pub fn compute_topological_signature(
    graph: &PersistentGraph,
    config: &TopologyConfig,
) -> TopologicalSignature;

pub fn dtl_evaluate(
    graph: &PersistentGraph,
    spectral: &SpectralDecomposition,
) -> BTreeMap<String, bool>;

pub fn dedup_filter(
    observations: &[Observation],
    seen: &mut BTreeSet<Hash256>,
) -> (Vec<Observation>, usize);  // (unique, dup_count)
\end{verbatim}


\section{Acceptance Tests}\label{sec:topo:tests}

\begin{longtable}{l l p{7.5cm}}
\toprule
\textbf{ID} & \textbf{Name} & \textbf{Procedure} \\
\midrule
\endfirsthead
\midrule
\endhead
\bottomrule
\endfoot
AT-T1 & Laplacian correctness & Build 5-vertex graph; verify $L$
  has zero row sums and correct eigenvalue count. \\
AT-T2 & Spectral gap connected & Build connected graph; verify
  $\Delta > 0$. Build disconnected graph; verify $\Delta = 0$. \\
AT-T3 & CTQW self-return & Initialize CTQW at vertex $j$ at $t=0$;
  verify $p_{jj}(0) = 1$. \\
AT-T4 & CTQW unitarity & Verify $\sum_k p_{jk}(t) = 1$ for all $j$
  and sampled $t$ values (within floating-point tolerance). \\
AT-T5 & Kuramoto convergence & Initialize identical phases; verify
  $r = 1$. Initialize random phases with strong coupling; verify
  $r \to 1$ over sufficient steps. \\
AT-T6 & Kuramoto incoherence & Initialize uniformly distributed phases
  with zero coupling ($\kappa = 0$); verify $r \approx 0$. \\
AT-T7 & DTL connected predicate & Build connected and disconnected
  graphs; verify DTL predicate matches spectral gap. \\
AT-T8 & Fixpoint detection & Feed identical vertex sets for
  $n_{\mathrm{consec}}$ iterations; verify \texttt{Converged}. Feed
  changing sets; verify \texttt{NotYet}. \\
AT-T9 & Deduplication & Ingest 100 observations with 20 duplicates;
  verify 80 pass through, 20 logged as duplicates. \\
AT-T10 & Topological signature determinism & Compute signature twice
  on same graph; verify bitwise identical. \\
AT-T11 & Budget fallback & Set budget to 1ms on a large graph; verify
  fallback is used and logged. \\
AT-T12 & Crystal signature enrichment & Run S-Basic with topology
  enabled; verify crystals contain spectral gap, Kuramoto $r$, and
  mean propagation time in $\tau(C)$. \\
\end{longtable}


%%% =====================================================================
\part{Crate C17: \texttt{isls-store}}\label{part:store}

\section{Purpose}\label{sec:store:purpose}

The store crate upgrades \ISLS{} from filesystem-based persistence
(\texttt{\~{}/.isls/} with JSONL files) to a structured SQLite database.
This enables: project management, run indexing, crystal querying,
settings persistence, and maintenance operations. It is the foundation
for any future GUI, REST API, or multi-user deployment.

\section{Data Model}\label{sec:store:model}

\begin{definition}[Store Entities]\label{def:storeentities}
The store manages the following entity types:
\begin{itemize}[nosep]
  \item \textbf{Project}: a named workspace containing runs, crystals,
        and configuration.
  \item \textbf{Run}: a single execution (discover or execute mode)
        with associated manifest.
  \item \textbf{Crystal}: a committed semantic crystal with all fields
        from the crystal schema.
  \item \textbf{Trace}: per-tick trace entries for a run.
  \item \textbf{Manifest}: execution manifest for a run.
  \item \textbf{Capsule}: sealed capsule objects.
  \item \textbf{Registry Snapshot}: point-in-time copy of all
        registries used in a run.
  \item \textbf{Metric}: time-series metric snapshots.
  \item \textbf{Alert}: emitted alert events.
  \item \textbf{Setting}: key-value configuration.
\end{itemize}
\end{definition}

\section{Schema}\label{sec:store:schema}

\begin{verbatim}
-- isls-store schema v1

CREATE TABLE projects (
    id          TEXT PRIMARY KEY,   -- content-addressed
    name        TEXT NOT NULL,
    created_at  TEXT NOT NULL,      -- ISO 8601
    description TEXT DEFAULT ''
);

CREATE TABLE runs (
    id            TEXT PRIMARY KEY,   -- manifest.run_id
    project_id    TEXT NOT NULL REFERENCES projects(id),
    mode          TEXT NOT NULL,      -- 'discover' | 'execute'
    rd_digest     TEXT NOT NULL,
    started_at    TEXT NOT NULL,
    finished_at   TEXT,
    ticks         INTEGER NOT NULL,
    crystal_count INTEGER DEFAULT 0,
    status        TEXT DEFAULT 'running' -- 'running'|'completed'|'failed'
);

CREATE TABLE crystals (
    crystal_id         TEXT PRIMARY KEY,
    run_id             TEXT NOT NULL REFERENCES runs(id),
    stability_score    REAL NOT NULL,
    free_energy        REAL NOT NULL,
    created_at_tick    INTEGER NOT NULL,
    carrier_instance   INTEGER NOT NULL,
    constraint_count   INTEGER NOT NULL,
    region_size        INTEGER NOT NULL,
    topology_signature TEXT NOT NULL,    -- JSON
    validation_status  TEXT DEFAULT 'pending', -- 'pending'|'passed'|'failed'
    data               TEXT NOT NULL     -- full crystal JSON
);

CREATE TABLE traces (
    id              INTEGER PRIMARY KEY AUTOINCREMENT,
    run_id          TEXT NOT NULL REFERENCES runs(id),
    tick            INTEGER NOT NULL,
    input_digest    TEXT NOT NULL,
    state_digest    TEXT NOT NULL,
    crystal_id      TEXT,               -- NULL if no crystal this tick
    gate_snapshot   TEXT NOT NULL,       -- JSON
    metrics_digest  TEXT NOT NULL,
    UNIQUE(run_id, tick)
);

CREATE TABLE manifests (
    run_id           TEXT PRIMARY KEY REFERENCES runs(id),
    manifest_json    TEXT NOT NULL,      -- full manifest JSON
    verification     TEXT DEFAULT 'pending'
);

CREATE TABLE capsules (
    id              TEXT PRIMARY KEY,    -- content-addressed
    run_id          TEXT REFERENCES runs(id),
    policy_json     TEXT NOT NULL,
    created_at      TEXT NOT NULL,
    opened_count    INTEGER DEFAULT 0,
    max_uses        INTEGER,
    expires_at      TEXT
);

CREATE TABLE registries (
    id              TEXT PRIMARY KEY,
    run_id          TEXT NOT NULL REFERENCES runs(id),
    kind            TEXT NOT NULL,       -- 'operator'|'profile'|...
    digest          TEXT NOT NULL,
    data            TEXT NOT NULL        -- full registry JSON
);

CREATE TABLE metrics (
    id              INTEGER PRIMARY KEY AUTOINCREMENT,
    run_id          TEXT NOT NULL REFERENCES runs(id),
    tick            INTEGER NOT NULL,
    snapshot_json   TEXT NOT NULL,       -- full MetricSnapshot JSON
    UNIQUE(run_id, tick)
);

CREATE TABLE alerts (
    id              INTEGER PRIMARY KEY AUTOINCREMENT,
    run_id          TEXT NOT NULL REFERENCES runs(id),
    tick            INTEGER NOT NULL,
    metric_id       TEXT NOT NULL,
    level           TEXT NOT NULL,       -- 'green'|'yellow'|'red'
    message         TEXT NOT NULL
);

CREATE TABLE settings (
    key             TEXT PRIMARY KEY,
    value           TEXT NOT NULL,
    updated_at      TEXT NOT NULL
);
\end{verbatim}


\section{Invariant Preservation}\label{sec:store:invariants}

\begin{invariant}[Store Append-Only Crystals]\label{inv:storeappend}
The crystals table \textbf{MUST NOT} support \texttt{UPDATE} or
\texttt{DELETE} on the \texttt{crystal\_id}, \texttt{data},
\texttt{stability\_score}, \texttt{free\_energy}, or
\texttt{topology\_signature} columns. Only
\texttt{validation\_status} may be updated (from \texttt{pending} to
\texttt{passed} or \texttt{failed}). This preserves Invariant~I10
(Commit Immutability).
\end{invariant}

\begin{invariant}[Store Deterministic Queries]\label{inv:storedeterm}
All queries used in replay-relevant paths \textbf{MUST} include
explicit \texttt{ORDER BY} clauses on deterministic columns
(\texttt{crystal\_id}, \texttt{tick}, or \texttt{id}). No query may
depend on insertion order or SQLite internal rowid ordering for
correctness.
\end{invariant}

\begin{requirement}[Migration Framework]\label{req:migration}
The store \textbf{MUST} include a migration framework that:
\begin{enumerate}[nosep,label=(\alph*)]
  \item tracks applied migrations by version number,
  \item applies migrations in order at startup,
  \item is idempotent (re-running a migration on an already-migrated
        database is a no-op),
  \item preserves all existing data during schema evolution.
\end{enumerate}
\end{requirement}


\section{Rust API}\label{sec:store:api}

\begin{verbatim}
// isls-store/src/lib.rs

pub struct IslandStore {
    conn: Mutex<rusqlite::Connection>,
}

impl IslandStore {
    pub fn open(path: &Path) -> Result<Self>;
    pub fn open_memory() -> Result<Self>;  // for tests
    pub fn migrate(&self) -> Result<()>;

    // Projects
    pub fn create_project(&self, name: &str, desc: &str) -> Result<String>;
    pub fn list_projects(&self) -> Result<Vec<Project>>;
    pub fn get_project(&self, id: &str) -> Result<Project>;

    // Runs
    pub fn create_run(&self, project_id: &str, mode: &str,
                      rd_digest: &str, ticks: u64) -> Result<String>;
    pub fn finish_run(&self, run_id: &str, crystal_count: u64) -> Result<()>;
    pub fn get_run(&self, run_id: &str) -> Result<Run>;
    pub fn list_runs(&self, project_id: &str) -> Result<Vec<Run>>;

    // Crystals (append-only)
    pub fn insert_crystal(&self, crystal: &CrystalRow) -> Result<()>;
    pub fn get_crystal(&self, crystal_id: &str) -> Result<CrystalRow>;
    pub fn list_crystals(&self, run_id: &str) -> Result<Vec<CrystalRow>>;
    pub fn update_validation(&self, crystal_id: &str,
                             status: &str) -> Result<()>;

    // Traces
    pub fn insert_trace(&self, trace: &TraceRow) -> Result<()>;
    pub fn get_traces(&self, run_id: &str) -> Result<Vec<TraceRow>>;

    // Manifests
    pub fn insert_manifest(&self, run_id: &str,
                           json: &str) -> Result<()>;
    pub fn get_manifest(&self, run_id: &str) -> Result<String>;

    // Capsules
    pub fn insert_capsule(&self, capsule: &CapsuleRow) -> Result<()>;
    pub fn increment_opened(&self, capsule_id: &str) -> Result<()>;

    // Metrics
    pub fn insert_metric(&self, run_id: &str, tick: u64,
                         snapshot: &str) -> Result<()>;
    pub fn get_metrics(&self, run_id: &str) -> Result<Vec<MetricRow>>;
    pub fn get_latest_metric(&self, run_id: &str) -> Result<MetricRow>;

    // Alerts
    pub fn insert_alert(&self, alert: &AlertRow) -> Result<()>;
    pub fn get_alerts(&self, run_id: &str) -> Result<Vec<AlertRow>>;

    // Settings
    pub fn set_setting(&self, key: &str, value: &str) -> Result<()>;
    pub fn get_setting(&self, key: &str) -> Result<Option<String>>;

    // Maintenance
    pub fn vacuum(&self) -> Result<()>;
    pub fn integrity_check(&self) -> Result<bool>;
    pub fn export_run_zip(&self, run_id: &str,
                          path: &Path) -> Result<()>;
}
\end{verbatim}

\section{CLI Integration}\label{sec:store:cli}

\begin{verbatim}
# Initialize store (creates ~/.isls/isls.db)
isls init --store sqlite

# List projects
isls project list

# Create project
isls project create --name "market-analysis"

# Run within a project
isls run --mode shadow --ticks 100 --project market-analysis

# Query crystals
isls crystal list --run <run-id>
isls crystal show <crystal-id>

# Export a run as ZIP (manifest + crystals + traces + metrics)
isls export --run <run-id> --output run-export.zip

# Maintenance
isls store vacuum
isls store check
\end{verbatim}


\section{Acceptance Tests}\label{sec:store:tests}

\begin{longtable}{l l p{7.5cm}}
\toprule
\textbf{ID} & \textbf{Name} & \textbf{Procedure} \\
\midrule
\endfirsthead
\midrule
\endhead
\bottomrule
\endfoot
AT-D1 & Create and query project & Create project; list; verify
  present. \\
AT-D2 & Crystal append-only & Insert crystal; attempt UPDATE on data
  column; verify rejection. \\
AT-D3 & Run lifecycle & Create run; insert crystals; finish run; verify
  crystal\_count matches. \\
AT-D4 & Manifest round-trip & Insert manifest JSON; retrieve; verify
  identical. \\
AT-D5 & Trace ordering & Insert traces out of order; query; verify
  returned in tick order. \\
AT-D6 & Migration idempotent & Run migrate twice; verify no error and
  same schema. \\
AT-D7 & Integrity check & Populate database; run integrity\_check;
  verify true. \\
AT-D8 & Export ZIP & Run S-Basic into store; export ZIP; verify
  contains manifest.json, crystals.jsonl, traces.jsonl. \\
\end{longtable}


%%% =====================================================================
\part{Dependency Graph and Build Order}

\section{Updated Dependency Graph}

\begin{verbatim}
isls-types        <-- (all crates)
    ^
isls-registry     <-- (all crates that use operators)
    ^
isls-observe  <-- isls-persist <-- isls-topology (C16)
                       ^                ^
                       |          isls-extract (uses spectral features)
                       |                ^
               isls-consensus <-- isls-carrier
                       ^
               isls-archive <-- isls-morph
                       ^
               isls-manifest
                       ^
               isls-capsule
                       ^
               isls-scheduler (uses fixpoint from C16)
                       ^
               isls-engine (+ execute + fixpoint integration)
                       ^
               isls-store (C17: SQLite, depends on types + archive)
                       ^
               isls-harness (extended with topology tests)
                       ^
               isls-cli (extended with store commands)
\end{verbatim}

\section{Build Order}

\begin{enumerate}[nosep]
  \item \texttt{isls-topology} (C16): depends on \texttt{isls-types},
        \texttt{isls-persist}. New external dependency: \texttt{nalgebra}
        (already in workspace for linear algebra).
  \item \texttt{isls-store} (C17): depends on \texttt{isls-types},
        \texttt{isls-archive}. New external dependency:
        \texttt{rusqlite} with \texttt{bundled} feature.
  \item Engine integration: wire topology into \texttt{macro\_step}
        (spectral enrichment before extraction, Kuramoto before gate
        evaluation).
  \item CLI integration: add store commands, topology metrics in
        reports.
  \item Harness integration: add AT-T1 through AT-T12 and AT-D1
        through AT-D8.
\end{enumerate}

\section{External Dependencies}

\begin{tabular}{l l l}
\toprule
\textbf{Crate} & \textbf{Version} & \textbf{Purpose} \\
\midrule
\texttt{nalgebra} & 0.33 & Eigendecomposition, matrix exponential \\
\texttt{rusqlite} & 0.32 & SQLite database (bundled) \\
\bottomrule
\end{tabular}


%%% =====================================================================
\part{Implementation Handoff}

\section{Test Count}

\begin{tabular}{l r r}
\toprule
\textbf{Phase} & \textbf{New Tests} & \textbf{Cumulative} \\
\midrule
Core (C1--C9) & 130 & 130 \\
Extension Phase 1 (C10--C15) & 31 & 161 \\
Extension Phase 2 (C16--C17) & 20 & $\ge$ 181 \\
\bottomrule
\end{tabular}

\section{Verification Sequence}

\begin{verbatim}
# Build
cargo build --workspace --release 2>&1 | grep -c warning  -> 0

# Test
cargo test --workspace -- --test-threads=1  -> >= 181 passed

# S-Basic with topology enrichment
rm -rf ~/.isls
cargo run --bin isls --release -- init --store sqlite
cargo run --bin isls --release -- project create --name test
cargo run --bin isls --release -- ingest --adapter synthetic --scenario S-Basic
cargo run --bin isls --release -- run --mode shadow --ticks 100 --project test
cargo run --bin isls --release -- validate formal

# Verify topology in crystal signatures
# Crystals must contain spectral_gap, kuramoto_coherence,
# mean_propagation_time in topology_signature

# Export
cargo run --bin isls --release -- export --run latest --output test.zip
# Verify ZIP contains manifest.json, crystals.jsonl, traces.jsonl
\end{verbatim}

\section{Completion Criteria}

Phase 2 is complete when:
\begin{enumerate}[nosep]
  \item \texttt{isls-topology} compiles, all AT-T tests pass.
  \item \texttt{isls-store} compiles, all AT-D tests pass.
  \item Crystals from S-Basic contain hardened topological signatures
        with $\Delta$, $r_k$, $\bar{t}_{\mathrm{prop}}$.
  \item M25--M27 appear in metric snapshots and HTML reports.
  \item Store CLI commands work (init, project, crystal list, export).
  \item Total tests $\ge 181$, zero warnings.
  \item Full scenario script runs end-to-end with store backend.
\end{enumerate}


\vfill
\begin{center}
\rule{0.5\textwidth}{0.4pt}\\[0.5cm]
{\small Generated for \ISLS{} Extension Phase 2 v1.0.0.\\
Deterministic, append-only, replay-verified.}
\end{center}

\end{document}
